\scanpage{4}\mspage{2}
définissent
\[
\begin{aligned}
P&\longrightarrow\underline{\mathrm{Hom}}_{\mathrm{Hot}}(Q,R)\\
Q&\unsure{\text{---}|\longrightarrow}\underline{\mathrm{Hom}}_{\mathrm{Hot}}(P,R)
\end{aligned}
\]
d'où en passant aux $\pi_0$
\[
\begin{aligned}
\pi_0(P)&\longrightarrow\mathrm{Hom}_{\mathrm{Hot}}(Q,R)\\
\pi_0(Q)&\text{---}\ \mathrm{Hom}_{\mathrm{Hot}}(P,R)
\end{aligned}
\]
\deleted{\unsure{donc}}\unsure{en} d'autres termes,
\[
\begin{aligned}
\overset{\deleted{\unsure{\mathrm{Hom}(b}}}{\pi_0(b,c)}&\longrightarrow\mathrm{Hom}_{\mathrm{Hot}}\bigl(\mathrm{Ch}_{\infty}(a,b),\mathrm{Ch}_{\infty}(\unsure{a|b},c)\bigr)\\
\pi_0(a,b)&\longrightarrow\mathrm{Hom}_{\mathrm{Hot}}\bigl(\mathrm{Ch}_{\infty}(b,c),\mathrm{Ch}_{\infty}(a,c)\bigr)
\end{aligned}
\]
\deleted{\unsure{On}} \deleted{\unsure{On en a}} \deleted{\unsure{notera}, pour $\alpha\in\pi_0(b,c)$, par $\alpha*$ \unsure{l'hom.} \unsure{l'homom.} associé de $\mathrm{Ch}_{\infty}(a,b)$ dans $\mathrm{Ch}_{\infty}($} \deleted{Soit $M_{\pi}$ la \unsure{catégorie}} On en déduit en particulier pour $a,b\in M_0$
\[
\left\{\begin{aligned}
1/\ \pi_0(a,a)&\longrightarrow\mathrm{Hom}_{\mathrm{Hot}}\bigl(\mathrm{Ch}_{\infty}(a,b),\mathrm{Ch}_{\infty}(a,b)\bigr)\\
2/\ \pi_0(a,a)&\longrightarrow\mathrm{Hom}_{\deleted{\unsure{\mathrm{Hot}}}\mathrm{Hot}}\bigl(\mathrm{Ch}_{\infty}(b,a),\mathrm{Ch}_{\infty}(b,a)\bigr)
\end{aligned}\right.\tag*{$*)$}
\]
et un objet \inserted{$\alpha_a$ de $\pi_0(a,a)$} est unité bilatère si ses deux images \unsure{soint|soit|sont} les \deleted{foncteurs} \inserted{\unsure{hom.|homs.}} identiques \unsure{de|des} $\mathrm{Ch}_{\infty}(a,b)$ \unsure{resp\textsuperscript{t}} $\mathrm{Ch}_{\infty}(b,a)$. \unsure{Revenant|Passant} aux $\pi_0$ \deleted{\unsure{de}} de ceux-ci, \unsure{cela|cel.} implique que $\alpha$ est unité bilatère pour la composition \deleted{dans $M_{\pi}$ (\unsure{i.e.} \unsure{p.}} des $\pi_0(x,y)$, relativement à $b$. Si on \unsure{veut} que ce soit le cas pour \uline{\unsure{tous} $b$}, \unsure{cela|cel.} signifie \inserted{(prenant $b=a$)} que $\alpha_a$ est un objet unité bilatère pour la structure de comp. associative \inserted{(\unsure{donc} de monoïde)} \deleted{\unsure{dans}} $\pi_0(a,a)$, donc que $\pi_0(a,a)$ est un monoïde \uline{unitaire}, ayant $\alpha$ pour objet unité. L'élément $\alpha_a$ est donc déterminé de façon unique. \deleted{Mais} Et on voit que $\alpha_a$ existe
% archive sheet number: 3

